\section{Introduction}\label{sec:intro}

Foundation models are pretrained once on broad data and then adapted to many tasks~\cite{bommasani}.
When the pretraining is supervised, the label set is itself a design choice: the same examples can be divided into
few classes or many, and that division determines which distinctions the network is rewarded for
learning~\cite{mahajan}.
In image classification, relabeling a fixed set of images along a class hierarchy costs little once the network is
fine-tuned~\cite{huh}, but more when its features are used frozen~\cite{hong}.
Such hierarchies are human taxonomies, however, and it is seldom clear what information a given merge removes from
the labels.

Particle physics allows the question to be put more sharply, because there the labels come from simulation.
At the Large Hadron Collider, quarks, gluons and the decay products of heavy particles reach the detector as jets,
collimated sprays of tens of particles, and identifying what produced a jet, known as jet tagging, is one of the most
developed uses of machine learning in the field~\cite{larkoski,feickert}, with graph networks~\cite{particlenet} and
transformers~\cite{part} on the jet's particles as the leading architectures.
Each simulated jet comes with an exact record of the process that produced it, so its label can be made as fine as
the simulated physics allows.
Networks pretrained on large sets of simulated jets for reuse in other tasks use self-supervised~\cite{jetclr,rs3l,mpm,hepjepa},
generative~\cite{omnijet} and supervised or combined~\cite{part,omnilearn,omnilearned} objectives.
Sophon~\cite{sophon}, a Particle Transformer pretrained by classification, uses the freedom of simulation: it sorts
the jets of the JetClass-II dataset~\cite{jetclasstwo} into \Kfull{} classes, which specify the particle that decayed
and what it decayed into.

Whether such fine labels make the pretrained model better has not been tested for jets in a controlled way.
In the comparisons published so far the training jets change together with the labels, so that a model trained on
fewer classes has also seen fewer kinds of jets~\cite{sophon,glopartdp,datacomp,superjepa}, and the two effects cannot
be told apart.
The answer decides how much generator-level detail future pretraining datasets must record as labels, and whether
analyses can read signals that are not among the classes from frozen features.

We pretrain Sophon's network on the same JetClass-II jets with four nested label sets of \Kfull, \Kres, \Kmid{} and
\Kcoarse{} classes (table~\ref{tab:sets}).
In the coarsest, for example, a decay to two bottom quarks and a decay to two charm quarks both become a decay to two
quarks or gluons.
We also pretrain on random groupings into \Kcoarse{} classes of the same sizes, which cut across this hierarchy, on
groupings that differ from the \Kcoarse-class set in a single flavor boundary, and with an auxiliary regression of the
jet mass.
Each nested label set is pretrained \VRunsWord{} times from different seeds.

We use each pretrained model in two ways.
With the network frozen, we train a linear classifier on its features (a linear probe)~\cite{alain}, for the
\Kfull{} pretraining classes and for pairs of classes that differ in a single flavor, and a small classifier for the
benchmark with which Sophon was introduced, the tagging of a hypothetical heavy particle $X$ that decays to
$b\bar b$ or to $bs$.
With the network fine-tuned, we measure the accuracy on the pretraining classes, on the ten classes of
JetClass~\cite{part}, which come from a different simulation, and on two public binary benchmarks, the identification
of jets from top quarks (top tagging)~\cite{toptag} and the discrimination of quark jets from gluon
jets~\cite{efn,qgpythia,qgherwig}.

The frozen features lose much of what the labels merge: classes that share a coarse label are identified less
accurately, pairs of such classes are separated less well, and a resonance search loses signal efficiency.
How much is lost depends more on which classes are merged than on how many, and a single measure of the confusable
classes a label set keeps apart tracks the accuracy of the label sets we study.
Fine-tuning on a limited labeled sample recovers most of the loss but not all of it.
After fine-tuning, the binary benchmarks, whose inputs carry little flavor information, show no dependence on the
label set, and the mass regression sharpens the mass resolution of the frozen features at a small cost in fine-tuned accuracy.
